\begin{definition}
\label{definition-regular-ideal-sheaf}
\begin{reference}
The concept of a Koszul-regular ideal sheaf was introduced in
\cite[Expose VII, Definition 1.4]{SGA6} where it was called a
regular ideal sheaf.
\end{reference}
Let $X$ be a ringed space. Let $\mathcal{J} \subset \mathcal{O}_X$
be a sheaf of ideals.
\begin{enumerate}
\item We say $\mathcal{J}$ is {\it regular} if for every
$x \in \text{Supp}(\mathcal{O}_X/\mathcal{J})$ there exists an open
neighbourhood $x \in U \subset X$ and a regular sequence
$f_1, \ldots, f_r \in \mathcal{O}_X(U)$ such that $\mathcal{J}|_U$
is generated by $f_1, \ldots, f_r$.
\item We say $\mathcal{J}$ is {\it Koszul-regular} if for every
$x \in \text{Supp}(\mathcal{O}_X/\mathcal{J})$ there exists an open
neighbourhood $x \in U \subset X$ and a Koszul-regular sequence
$f_1, \ldots, f_r \in \mathcal{O}_X(U)$ such that $\mathcal{J}|_U$
is generated by $f_1, \ldots, f_r$.
\item We say $\mathcal{J}$ is {\it $H_1$-regular} if for every
$x \in \text{Supp}(\mathcal{O}_X/\mathcal{J})$ there exists an open
neighbourhood $x \in U \subset X$ and a $H_1$-regular sequence
$f_1, \ldots, f_r \in \mathcal{O}_X(U)$ such that $\mathcal{J}|_U$
is generated by $f_1, \ldots, f_r$.
\item We say $\mathcal{J}$ is {\it quasi-regular} if for every
$x \in \text{Supp}(\mathcal{O}_X/\mathcal{J})$ there exists an open
neighbourhood $x \in U \subset X$ and a quasi-regular sequence
$f_1, \ldots, f_r \in \mathcal{O}_X(U)$ such that $\mathcal{J}|_U$
is generated by $f_1, \ldots, f_r$.
\end{enumerate}
\end{definition}